In this section, we present proof certificates and automata-based methods for the verification of LTL properties.
% We first introduce two types of subsets that are used through out the work, which match sets in the Büchi acceptance condition (Section 2.3).

% \begin{definition}
% ($\mathcal{X}_u$ and $\mathcal{X}_{VF}$).  A subset $\mathcal{X}_{u} \subseteq \mathcal{X}$ is an unsafet set just if elements of $\mathcal{X}_{u}$ can not occur in any trajectory of $\mathfrak{S}$. A subset $\mathcal{X}_{VF} \subseteq \mathcal{X}$ is a finitely visited set if elements of $\mathcal{X}_{VF}$ must occur only finitely often in any trajectory of $\mathfrak{S}$.
% \end{definition}



\subsection{State Persistence Certificate}

\begin{definition}[state persistence certificate]
  \label{systempc}
  Consider a system $\mathfrak{S} = (\mathcal{X}, \mathcal{X}_0, f)$.A bounded function $B_{sp}(x): \mathcal{X} \rightarrow \mathbb{R}_{\geq 0}$ is a \emph{state persistence certificate} for $\mathfrak{S}$ with respect to $\mathcal{X}_{VF} \subseteq \mathcal{X}$ if there exists a value $\epsilon > 0$ such that for $x \in \mathcal{X}, x_0 \in \mathcal{X}_0, x^{\prime} = f(x)$, we have:\\
  \begin{align}
    &B_{sp}(x_0) \geq 0, \label{bspnn} \\
    &B_{sp}(x) \geq 0 \Rightarrow B_{sp}(x^{\prime}) \geq 0, \label{bsptnn} \\
    &B_{sp}(x) \geq B_{sp}(x^{\prime}) + I_{\mathcal{X}_{VF}}(x^{\prime})\epsilon, \label{bspni1}\\
    &B_{sp}(x) \geq B_{sp}(x^{\prime}) + I_{\mathcal{X}_{VF}}(x)\epsilon. \label{bspni2}
  \end{align}
  We have two types of \emph{the state persistence certificate} based on the different conditions, but both types can be used to prove persistence.
  If $B_{sp}$ satisfies (\ref{bspnn}, \ref{bsptnn}, \ref{bspni1}), then $B_{sp}$ is a \emph{$\mathrm{type~1}$ state persistence certificate}. If $B_{sp}$ satisfies (\ref{bspnn}, \ref{bsptnn}, \ref{bspni2}), then $B_{sp}$ is a \emph{$\mathrm{type~2}$ state persistence certificate}.
\end{definition}

\begin{theorem}[persistence]
  Given a system $\mathfrak{S}=(\mathcal{X}, \mathcal{X}_0, f)$ and $\mathcal{X}_{VF} \subseteq \mathcal{X}$, if a $\mathrm{type~1}$ state persistence certificate $B_{sp}$ with respect to $\mathcal{X}_{VF}$ can be found, then it can be guaranteed that $\tau$ visits(\ref{notation123}) \(\mathcal{X}_{VF}\) finitely often for any \emph{trajectory} $\tau = \langle x_0, x_1, \dots \rangle \in \mathfrak{T}(\mathfrak{S})$.
\end{theorem}

\begin{proof}
We assume there exists a trajectory $\tau$ that visits $\mathcal{X}_{VF}$ infinitely often. We extract the states in $\tau$ that visit $\mathcal{X}_{VF}$ in order to construct $\tau^{\prime} = \langle y_0, y_1, \dots, y_n, \dots \rangle$, where $y_i \in \mathcal{X}_{VF}$. We now prove that $B_{sp}(y_i) \leq B_{sp}(x_0) - i\epsilon$ for all $i \in \mathbb{N}$.
We proceed by induction on $i$.
\begin{itemize}
    \item \textbf{Base case:} When $i = 0$, we have $B_{sp}(y_0) \leq B_{sp}(x_0)$ according to the Lemma (\ref{state-persistence-lemma-b0-top}).
    \item \textbf{Inductive step:} We assume the statement holds for $i = k$, i.e., $B_{sp}(y_k) \leq B_{sp}(x_0) - k\epsilon$. When $i = k + 1$:
    \begin{itemize}
        \item If $y_{k+1} = f(y_k)$, we have $B_{sp}(y_{k+1}) \leq B_{sp}(y_k) - \epsilon$ according to the property (\ref{bspni1}). Combining this with the inductive hypothesis, we obtain $B_{sp}(y_{k+1}) \leq B_{sp}(x_0) - (k + 1)\epsilon$.
        \item If $y_{k+1} \neq  f(y_k)$, i.e., $y_k$ is not the \emph{predecessor}(\ref{dtds}) of $y_{k+1}$, so $y_{k+1}^{pre}$ must appear after $y_k$ in $\tau$.
        So we have $B_{sp}(y_{k+1}^{pre}) \leq B_{sp}(y_k)$ according to the Lemma (\ref{state-peristence-lemma-non-increase}).
        We also have $B_{sp}(y_{k+1}^{pre}) \geq  B_{sp}(y_{k+1}) + \epsilon$ according to the property (\ref{bspni1}).
        Combining these, we obtain $B_{sp}(y_{k+1}) + \epsilon \leq B_{sp}(y_k)$.
        Using the inductive hypothesis, we have $B_{sp}(y_{k+1}) \leq B_{sp}(x_0) - (k + 1)\epsilon$.
    \end{itemize}
\end{itemize}
Consequently, there exists $n \in \mathbb{N}$ such that $B_{sp}(y_n) \leq B_{sp}(x_0) - n\epsilon < 0$. This contradicts the Lemma (\ref{bspnnlemma0}).
\end{proof}

We have shown that \emph{$\mathrm{type~1}$ state persistence certificate} is sufficient to guarantee \emph{persistence}(\ref{specification123}). By using a similar proof method, we can easily prove that the \emph{$\mathrm{type~2}$ state persistence certificate} is able to achieve the same conclusion. That is, both types of state persistence certificates can ensure \emph{persistence}. Next, we will demonstrate that the \emph{state persistence certificate} possesses stronger expressive power compared with the closure certificate for persistence.

\subsection{State Safety Certificate}

\begin{definition}[state safety certificate]
  \label{state-safety-certificate-definition}
  Given a system $\mathfrak{S} = (\mathcal{X}, \mathcal{X}_0, f)$, a labelling function $\mathfrak{L}:\mathcal{X} \rightarrow \Sigma$ and an NBA $\mathcal{A}=(\Sigma, Q, q_0, \delta, Acc)$. Construct the product system $\mathfrak{S} \times \mathcal{A} = (\mathcal{Y}, \mathcal{Y}_0, F)$, where $\mathcal{Y} = \{(x, q) \,|\, x \in \mathcal{X}, q \in Q \}$, $\mathcal{Y}_0 = \{ (x_0, q_0) \,|\, x_0 \in \mathcal{X}_0 \}$, and $F(x,q) = \{(x^{\prime}, q') \,|\, x^{\prime} = f(x) \land (q, \mathfrak{L}(x), q') \in \delta\}$. A bounded function $B_{k}(x, q): \mathcal{X} \times Q \rightarrow \mathbb{R}$ is a \emph{state safety certificate} with respect to $q_k$, where $(x,q) \in \mathcal{Y}, (x_0, q_0) \in \mathcal{Y}_0, q_k \in Q, \mathcal{Y}_u = \{(x, q_k) \,|\, x \in \mathcal{X}\}, (x_0, q_0) \in \mathcal{Y}_0, (x_u, q_u) \in \mathcal{Y}_u, (x^{\prime}, q^{\prime}) \in F(x, q)$, we have:
  \begin{align}
    &B_{k}(x_0, q_0) \geq 0, \label{bsinn1} \\ % initial state is non-negative
    &B_{k}(x_u, q_u) < 0, \label{bsun1} \\ % accepting state is negative
    &B_{k}(x, q) \geq 0 \Rightarrow B_{k}(x^{\prime}, q^{\prime}) \geq 0. \label{bstnn1} % transition will remain non-negative
  \end{align}
\end{definition}

\subsection{Transition Persistence Certificate}
\begin{definition}[transition persistence certificate]
  \label{bsbstct}
  Consider a system $\mathfrak{S} = (\mathcal{X}, \mathcal{X}_0, f)$, a labelling function $\mathfrak{L}:\mathcal{X} \rightarrow \Sigma$ and an NBA $\mathcal{A}=(\Sigma, Q, q_0, \delta, Acc)$, we define $\Delta_{k, l} = \{(q_k, \sigma, q_l) \mid \mathcal{P}(q_k, \sigma, q_l) \land (q_k, \sigma, q_l) \in \delta\}$
  as the set of transitions that satisfy the predicate $\mathcal{P}$ and move from $q_k$ to $q_l$.
  A bounded function $B_{\Delta_{k, l}}(x): \mathcal{X} \rightarrow \mathbb{R}$ is a \emph{transition persistence certificate} with
  respect to $\Delta_{k,l}$ where $x_0 \in \mathcal{X}_0, x \in \mathcal{X}, x^{\prime} = f(x)$, we have:
  \begin{align}
    &B_{\Delta_{k, l}}(x_0) \geq 0, \label{bstnn} \\
    &B_{\Delta_{k, l}}(x) \geq 0 \Rightarrow B_{\Delta_{k, l}}(x^{\prime}) \geq 0, \label{bsttnn} \\
    &B_{\Delta_{k, l}}(x) \geq B_{\Delta_{k, l}}(x^{\prime}) + I_{\Delta_{k,l}}(q_k, \mathfrak{L}(x), q_l)\epsilon. \label{bstni}
  \end{align}
\end{definition}

\subsection{Discussion}
\begin{theorem}[LTL VERIFICATION]
  \label{ltl-thm}
   Given a system $\mathfrak{S} = (\mathcal{X}, \mathcal{X}_0, f)$, a labelling function $\mathfrak{L}: \mathcal{X} \rightarrow \Sigma$ and an NBA $\mathcal{A} = (\Sigma, Q, q_0, \delta, Acc)$ representing the negation of an LTL formula $\phi$. For each accepting transition start state $q_k \in Acc^{start}$ (\ref{atstart}), if we can find either:
  \begin{enumerate}
      \item A corresponding state safety certificate $B_k$, or
      \item Transition persistence certificates $B_{\Delta_{k,l}}$ for all $(k, l) \in Acc^{pair}$ (\ref{atpair}).
  \end{enumerate}
  then it can be guaranteed that $\mathfrak{S} \models_{\mathfrak{L}} \phi$.
\end{theorem}
\begin{proof}
  Given an accepting transition start state $q_k \in Acc^{start}$:
  \begin{enumerate}
      \item If a state safety certificate $B_k$ is found, by Lemma~(\ref{state-safety-lemma}), all runs of words in $\mathfrak{L}(\mathfrak{S})$ will never visit state $q_k$, preventing all accepting transitions from $q_k$.

      \item For each $(k, l) \in Acc^{pair}$, if transition persistence certificate $B_{\Delta_{k,l}}$ is found where $\Delta_{k,l} = Acc^{k,l}$ (\ref{atkl}), by Lemma~(\ref{transition-persistence-lemma1}), all runs will make accepting transitions from $q_k$ to $q_l$ happen finitely often.
  \end{enumerate}

  This ensures that all accepting transitions in $Acc$ occur only finitely often. Therefore, the system language is rejected by automaton $\mathcal{A}$, implying $\mathfrak{L}(\mathfrak{S}) \cap \mathfrak{L}(\mathcal{A}) = \emptyset$. Since $\mathfrak{L}(\mathcal{A}) = \mathfrak{L}(\neg \phi)$, we have $\mathfrak{L}(\mathfrak{S}) \subseteq \mathfrak{L}(\phi)$, and consequently $\mathfrak{S} \models_{\mathfrak{L}} \phi$.
\end{proof}

Theorem~(\ref{ltl-thm}) illustrates the verification of LTL properties using state safety certificates and transition persistence certificates. By constructing the product system $\mathfrak{S} \times \mathcal{A}$ to synchronize the states of system $\mathfrak{S}$ and automaton $\mathcal{A}$, the state safety certificate $B_k$ establishes a barrier over this product system to block the automaton from reaching state $q_k$. The transition persistence certificate $B_{\Delta_{k,l}}$ demonstrates that accepting transitions are visited only finitely many times for any trajectory.

%  For LTL property verification, the efficiency can be further improved using transition-based NBA.
Both state safety certificates and transition persistence certificates can be utilized to prove LTL properties. To illustrate the necessity of the two types of certificates, consider the following scenario where the synthesis of a transition persistence certificate is not feasible. Assume the system indeed satisfies the LTL specification $G \lnot a$, but the transition persistence certificate cannot be synthesized. To prove the property, we firstly construct an NBA $\mathcal{A}$ for $F a$. Note that in Figure \ref{pic-tbnba1}, the alphabet is $\{ \{\},\{a\}\}$, the accepting transitions are $(q_0, \{\}, q_0)$ and $(q_0, \{a\}, q_0)$, so any letter in alphabet as input at $q_0$ only transitions to $q_0$. According to the properties (\ref{bstnn}, \ref{bsttnn}, \ref{bstni}) of transition persistence certificate, $B_{0,0}(x_0) \geq 0$, $B_{0,0}(x) \geq 0 \Rightarrow B_{0,0}(x') \geq 0$, and $B_{0,0}(x) \geq B_{0,0}(x') + \epsilon$ should be satisfied. In particular, the third condition can lead to an infinite descent of the certificate, ultimately resulting in a contradiction with Lemmas (\ref{bspnnlemma0}, \ref{state-persistence-lemma-b0-top}) and making it impossible to synthesize the certificate. This problem is caused by reachability: if the \emph{start state} $q_0$ of the accepting transitions cannot be reached by all runs. However, this problem can be solved using the state safety certificate. By determining whether state $q_0$ can be reached from state $q_1$, if a state safety certificate can be synthesized, it guarantees that $q_0$ is never reached, implying that it is finitely visited. If there exists a \emph{trajectory} that allows $\mathcal{A}$ to visit $q_0$, then the state safety certificate cannot be synthesized. In such cases, the transition persistence certificate can serve as a witness that the relevant accepting transitions are finitely visited. Therefore, the combination of these two can complement each other's limitations.
